%%%%%%%%%%  scan idx 210  |  folio 197  %%%%%%%%%%
donc que l'espace sous-jacent à $X$ est noethérien de dimension $2$. Soit $I$ un $A$-module injectif tel que $H^2_Y(X,\widetilde I)\ne0$, et notons $U$ l'ouvert complémentaire de $Y$ dans $X$. On a (SGA 2 II 4)
\[
 H^2_Y(X,I)\simeq H^1(U,I)\ne0.
\]
Or, comme $U$ est séparé, le foncteur $\mathrm D^+(\Qcoh(U))\to\mathrm D^+(U)$ est pleinement fidèle (II 3.5), donc $\widetilde I|U$ n'est pas injectif dans $\Qcoh(U)$. On a donc prouvé l'assertion (0.1).

Notons $Z$ le schéma quasi-compact et quasi-séparé obtenu en recollant deux exemplaires de $X$ le long de $U$, $j_1,j_2$ les deux immersions ouvertes canoniques de $X$ dans $Z$, de sorte qu'on a un carré cartésien
\[
\begin{tikzcd}
 U \arrow[r,"i"] \arrow[d,"i"'] & X \arrow[d,"j_1"]\\
 X \arrow[r,"j_2"] & Z .
\end{tikzcd}
\]
On a
\[
 j_2^*R^1j_{1*}(\widetilde I)\simeq R^1i_*(i^*\widetilde I),
\]
d'où
\[
 H^0(X,R^1j_{1*}(\widetilde I))\simeq H^1(U,\widetilde I)\ne0,
\]
donc $R^1j_{1*}(\widetilde I)\ne0$, ce qui prouve (0.2).

D'autre part, on a, d'après (II 3.4 et 3.5)
\[
 \mathbb RQ\,\mathbb Rj_{1*}(\widetilde I)\simeq j_{1*}(\widetilde I).
\]
On en déduit une suite spectrale, dont la suite exacte des termes de bas degrés fournit l'isomorphisme
\[
 R^1j_{1*}(\widetilde I)\simeq R^2Q(j_{1*}\widetilde I).
\]
Le faisceau $j_{1*}\widetilde I$ étant injectif dans $\Qcoh(Z)$, on a prouvé la première 
par%
